% ========================================================
% BÀI TẬP TOÁN 5 - KẾT NỐI TRI THỨC VỚI CUỘC SỐNG
% Bài 56: Các đơn vị đo thời gian
% Định dạng: BTTeX Part III - Trắc nghiệm điền số ngắn
% ========================================================

% Câu 1
\begin{ex}
  Một thế kỉ bằng bao nhiêu năm?
  \shortans{100}
  \loigiai{
    1 thế kỉ = 100 năm.
  }
\end{ex}

% Câu 2
\begin{ex}
  Năm 2024 thuộc thế kỉ thứ bao nhiêu?
  \shortans{21}
  \loigiai{
    Năm 2024 thuộc thế kỉ XXI (21).
  }
\end{ex}

% Câu 3
\begin{ex}
  Năm nhuận có bao nhiêu ngày?
  \shortans{366}
  \loigiai{
    Năm nhuận có 366 ngày (tháng 2 có 29 ngày).
  }
\end{ex}

% Câu 4
\begin{ex}
  Năm không nhuận có bao nhiêu ngày?
  \shortans{365}
  \loigiai{
    Năm thông thường có 365 ngày.
  }
\end{ex}

% Câu 5
\begin{ex}
  Tháng 7 có bao nhiêu ngày?
  \shortans{31}
  \loigiai{
    Tháng 7 có 31 ngày.
  }
\end{ex}

% Câu 6
\begin{ex}
  Tháng 4 có bao nhiêu ngày?
  \shortans{30}
  \loigiai{
    Tháng 4 có 30 ngày.
  }
\end{ex}

% Câu 7
\begin{ex}
  Một tuần lễ có bao nhiêu giờ?
  \shortans{168}
  \loigiai{
    $7 \times 24 = 168$ giờ.
  }
\end{ex}

% Câu 8
\begin{ex}
  Một ngày có bao nhiêu phút?
  \shortans{1440}
  \loigiai{
    $24 \times 60 = 1\,440$ phút.
  }
\end{ex}

% Câu 9
\begin{ex}
  Một giờ có bao nhiêu giây?
  \shortans{3600}
  \loigiai{
    $60 \times 60 = 3\,600$ giây.
  }
\end{ex}

% Câu 10
\begin{ex}
  Đổi $3$ ngày ra giờ được bao nhiêu giờ?
  \shortans{72}
  \loigiai{
    $3 \times 24 = 72$ giờ.
  }
\end{ex}

% Câu 11
\begin{ex}
  Đổi $2{,}5$ giờ ra phút được bao nhiêu phút?
  \shortans{150}
  \loigiai{
    $2{,}5 \times 60 = 150$ phút.
  }
\end{ex}

% Câu 12
\begin{ex}
  Đổi $0{,}75$ giờ ra phút được bao nhiêu phút?
  \shortans{45}
  \loigiai{
    $0{,}75 \times 60 = 45$ phút.
  }
\end{ex}

% Câu 13
\begin{ex}
  Đổi $1{,}2$ phút ra giây được bao nhiêu giây?
  \shortans{72}
  \loigiai{
    $1{,}2 \times 60 = 72$ giây.
  }
\end{ex}

% Câu 14
\begin{ex}
  Đổi $\frac{3}{4}$ giờ ra phút được bao nhiêu phút?
  \shortans{45}
  \loigiai{
    $\frac{3}{4} \times 60 = 45$ phút.
  }
\end{ex}

% Câu 15
\begin{ex}
  Đổi $\frac{1}{5}$ giờ ra phút được bao nhiêu phút?
  \shortans{12}
  \loigiai{
    $\frac{1}{5} \times 60 = 12$ phút.
  }
\end{ex}

% Câu 16
\begin{ex}
  Đổi $2$ thế kỉ rưỡi ($2{,}5$ thế kỉ) ra năm được bao nhiêu năm?
  \shortans{250}
  \loigiai{
    $2{,}5 \times 100 = 250$ năm.
  }
\end{ex}

% Câu 17
\begin{ex}
  Bác Hồ sinh năm 1890. Năm đó thuộc thế kỉ thứ bao nhiêu?
  \shortans{19}
  \loigiai{
    Năm 1890 thuộc thế kỉ XIX (19).
  }
\end{ex}

% Câu 18
\begin{ex}
  Chiến thắng Điện Biên Phủ năm 1954 thuộc thế kỉ thứ bao nhiêu?
  \shortans{19}
  \loigiai{
    Năm 1954 thuộc thế kỉ XIX (19).
  }
\end{ex}

% Câu 19
\begin{ex}
  Đổi $180$ phút ra giờ được bao nhiêu giờ?
  \shortans{3}
  \loigiai{
    $180 : 60 = 3$ giờ.
  }
\end{ex}

% Câu 20
\begin{ex}
  Đổi $300$ giây ra phút được bao nhiêu phút?
  \shortans{5}
  \loigiai{
    $300 : 60 = 5$ phút.
  }
\end{ex}

% Câu 21
\begin{ex}
  Đổi $36$ tháng ra năm được bao nhiêu năm?
  \shortans{3}
  \loigiai{
    $36 : 12 = 3$ năm.
  }
\end{ex}

% Câu 22
\begin{ex}
  Trong một năm dương lịch có bao nhiêu tháng có $31$ ngày?
  \shortans{7}
  \loigiai{
    Có 7 tháng: 1, 3, 5, 7, 8, 10, 12.
  }
\end{ex}

% Câu 23
\begin{ex}
  Trong một năm dương lịch có bao nhiêu tháng có đúng $30$ ngày?
  \shortans{4}
  \loigiai{
    Có 4 tháng: 4, 6, 9, 11.
  }
\end{ex}

% Câu 24
\begin{ex}
  Một tiết học kéo dài $40$ phút. Ba tiết học như thế kéo dài bao nhiêu phút?
  \shortans{120}
  \loigiai{
    $40 \times 3 = 120$ phút.
  }
\end{ex}

% Câu 25
\begin{ex}
  Đổi $1$ ngày $6$ giờ ra giờ được bao nhiêu giờ?
  \shortans{30}
  \loigiai{
    $24 + 6 = 30$ giờ.
  }
\end{ex}
